\[
x(t)=\frac{1+t^{10}}{1+t^2}
\]

تعریف \(T\) برای توان‌های صحیح \(t^k\) و خطی بودن روی ترکیب‌های جبری متناهی از آن‌ها اجازه می‌دهد ابتدا \(x(t)\) را به صورت جمع متناهی از دوجمله‌های \(t^k\) بنویسیم.

تجزیه‌ی جبری:
\[
1+t^{10}=(t^2+1)(t^8-t^6+t^4-t^2+1)
\]

بنابراین:
\[
x(t)=t^8-t^6+t^4-t^2+1
=\sum_{m=0}^{4}(-1)^m t^{2m}
\]

اگر ضرایب را صریح بنویسیم:
\[
x(t)=1\cdot t^0+(-1)t^2+1\cdot t^4+(-1)t^6+1\cdot t^8
\]

از قواعد \(T\{t^k\}=\cos(kt)\)، جمع‌پذیری و یکنواختی:
\[
T[x(t)]
=1\cdot\cos(0)+(-1)\cos(2t)+1\cdot\cos(4t)+(-1)\cos(6t)+1\cdot\cos(8t)
\]

چون \(\cos(0)=1\)، نتیجه نهایی:
\[
\boxed{
y(t)=T[x(t)]=1-\cos(2t)+\cos(4t)-\cos(6t)+\cos(8t)
}
\]
